\documentclass{article}
\usepackage{amsmath, amssymb, amsthm}
\usepackage{hyperref}
\usepackage{geometry}
\geometry{margin=1in}

\title{Erd\H{o}s Problem \#1038}
\date{Last updated: 2025-09-15}
\author{Source: erdosproblems.com}

\begin{document}
\maketitle

\noindent\textbf{Status:} OPEN \quad \textbf{No prize}

\noindent\textbf{Tags:} analysis

\medskip
\noindent\textbf{Problem Statement:}

Determine the infimum and supremum of\[\lvert \{ x\in \mathbb{R} : \lvert f(x)\rvert < 1\}\rvert\]as $f\in \mathbb{R}[x]$ ranges over all non-constant monic polynomials, all of whose roots are real and in the interval $[-1,1]$.

\bigskip
\noindent\textbf{Remarks:}

A problem of Erd\H{o}s, Herzog, and Piranian \cite{EHP58}, who proved that the measure of the set in question is always at most $2\sqrt{2}$ under the assumption that all the roots are in $\{-1,1\}$, and conjecture this is the best possible upper bound.

They also note that the infimum of the set in question is less than $2$, as witnessed by $f(x)=(x+1)(x-1)^m$ for $m\geq 3$. They further note that if the roots are restricted to $[-2,2]$ then the infimum is zero, as witnessed by a small perturbation of the Chebyshev polynomials. 

They further conjectured that, if the roots are restricted to $[-2,2]$, then\[\lvert \{ x\in \mathbb{R} : \lvert f(x)\rvert < 1\}\rvert\geq n^{-c}\]for an absolute constant $c>0$. This was proved by Pommerenke \cite{Po61}, who in fact showed that this set must contain an interval of width $\gg n^{-4}$.

The current best known bounds (see the discussion in the comments) are\[1.519\approx 2^{4/3}-1\leq \inf \leq 1.835\cdots\]and\[\sup = 2\sqrt{2}\approx 2.828.\]

\bigskip
\noindent\textbf{References:}

[EHP58] Erd\H{o}s, P. and Herzog, F. and Piranian, G., Metric properties of polynomials. J. Analyse Math. (1958), 125-148.
 
 [Po61] Pommerenke, Ch., On metric properties of complex polynomials. Michigan Math. J. (1961), 97-115.

\bigskip
\noindent\small{Source: \url{https://www.erdosproblems.com/1038}}
\end{document}
